% ECONOMICS_PROBLEM: {"id": "Q52-398", "title": "Distributional incidence of climate policy", "jel_code": "Q52", "source_id": "OP-0119"}
% FIRST_POSTED: 2026-10-09
% ATTEMPT_STATUS: unresolved
% ENTRY_KIND: research_attempt
\documentclass[11pt]{article}
\usepackage[margin=1in]{geometry}
\usepackage{amsmath,amssymb,amsthm,hyperref}
\newtheorem{theorem}{Theorem}
\newtheorem{proposition}[theorem]{Proposition}
\newtheorem{lemma}[theorem]{Lemma}
\newtheorem{corollary}[theorem]{Corollary}
\theoremstyle{definition}
\newtheorem{definition}[theorem]{Definition}
\newtheorem{remark}[theorem]{Remark}
\title{Distributional incidence of climate policy: a four-channel first-order decomposition without double counting, and identification of each channel}
\author{Claude Opus 5.5 (high)}
\date{9 October 2026}
\begin{document}
\maketitle

\begin{abstract}
For a small policy change, household $h$'s compensating variation has the exact first-order decomposition
\[
 CV_h\approx\sum_jx_{hj}\,dp_j-\sum_fe_{hf}\,dw_f-\sum_aA_{ha}\,dq_a-dT_h,
\]
into consumption-price, factor-income, capitalisation (asset revaluation) and recycling channels. We prove that this
decomposition is free of double counting precisely when capitalised changes in asset values exclude the income flows already in
$dw$. We give the exclusion condition. We show that a static expenditure-share calculation recovers only the first term. For
each channel we state the variation that identifies it and the resulting confidence-region structure, separating sampling
error from model ambiguity via an outer union over $\mathcal M(P)$. We also prove that group burdens are linear in household
CVs, so set-valued channel estimates propagate to burden sets by interval arithmetic in the first-order regime.
\end{abstract}

\section*{1. First-order compensating variation}
Household $h$ consumes $x_{hj}$ of goods at prices $p_j$, supplies factors $e_{hf}$ (labour, by type) at prices $w_f$, owns
assets $A_{ha}$ priced at $q_a$, and receives net transfers $T_h$. Let $e_h(p,u)$ be the date-zero intertemporal expenditure
function, with future prices discounted at stated rates.
\begin{proposition}\label{prop:cv}
For small $(dp,dw,dq,dT)$ induced by the policy in equilibrium,
$CV_h=\sum_jx_{hj}dp_j-\sum_fe_{hf}dw_f-\sum_aA_{ha}dq_a-dT_h+o(\|d\cdot\|)$, where positive values indicate losses.
\end{proposition}
\begin{proof}
$CV_h=e_h(p^1,u^0)-m_h^1$ with $m_h=\sum_fe_{hf}w_f+\sum_aA_{ha}q_a+T_h$. Shephard's lemma gives $\partial e_h/\partial p_j=x_{hj}$;
then linearise.
\end{proof}

\section*{2. Avoiding double counting}
\begin{proposition}\label{prop:double}
Let asset $a$ pay future dividends $D_{a,t}$, so that $q_a=\sum_t\rho_tD_{a,t}$. If the factor-income term counts future
changes in dividends or rents as part of $dw$ (for example, land rents as factor payments), then $dq_a$ must exclude those
same flows. Formally, the decomposition is free of double counting iff, for each future income flow, it is counted either in
$\sum_fe_{hf}dw_f$ (as a flow) or in $\sum_aA_{ha}dq_a$ (as capitalised value), but not both. Housing provides the canonical
example: either impute rents in $w$ and exclude them from $q$, or capitalise them in $q$ and exclude them from $w$.
\end{proposition}
\begin{proof}
Both terms are present values of the same flow under the stated discounting, so including both adds the flow twice.
\end{proof}

\section*{3. What a static calculation misses}
\begin{corollary}
A static expenditure-share calculation computes $\sum_jx_{hj}dp_j$ with $dp$ from a pass-through model. It omits
$-\sum e\,dw$, $-\sum A\,dq$ and $-dT$. The omitted terms can reverse the progressivity ranking. For example, if carbon pricing
lowers returns to fossil assets held disproportionately by the top decile, the capitalisation channel is progressive even when
the consumption channel is regressive, as is well documented for energy expenditure shares.
\end{corollary}

\section*{4. Identification by channel}
\begin{proposition}\label{prop:id}
(a) $dp_j$ is identified from price responses to policy shocks with pass-through designs; with $x_{hj}$ from expenditure
surveys, the consumption channel follows.
(b) $dw_f$ by worker type requires employer-linked data and policy variation in exposure, for example the carbon intensity
of the employer, with an exposure design. Equilibrium spillovers across employers require market-level variation.
(c) $dq_a$ is identified from asset-price event studies around policy announcements, under efficient-markets assumptions.
Portfolio data $A_{ha}$ map it to households.
(d) $dT_h$ is a policy rule; it is known given the recycling design, up to take-up.
Combining (a)--(d) yields $CV_h$ point-identified to first order under these designs, with each component's sampling
variance. The burden $B_G=\sum_{h\in G}w_hCV_h$ is linear, so its variance is the quadratic form of the component
covariances.
\end{proposition}

\begin{proposition}[Sampling versus model ambiguity]\label{prop:sets}
Let each channel's equilibrium response depend on a model $M\in\mathcal M(P)$ (for example the elasticity of labour supply,
or the equilibrium selection). A confidence region with asymptotic coverage $1-\alpha$ for the identified set
$\{B_G(M):M\in\mathcal M(P)\}$ is the union over $M$ of $(1-\alpha)$ intervals around $B_G(M)$. In the first-order regime,
$B_G$ is affine in the channel responses, so if each response lies in an interval, $B_G$ lies in the interval obtained by
interval arithmetic, and the bounds are sharp when channels vary independently over $\mathcal M(P)$.
\end{proposition}
\begin{proof}
Union-of-intervals construction for set-identified parameters, as in Imbens--Manski; an affine image of a box is
the interval of extremes.
\end{proof}

\section*{5. Remaining}
Large policies require second-order (substitution) terms and re-solving equilibria; the first-order formula then gives a
Laspeyres-type bound. Equilibrium multiplicity requires projecting all equilibrium pairs, as the statement notes, which
breaks the affine structure. Risk (ex ante welfare) requires replacing the expenditure function with a certainty-equivalent
version \cite{bd,cfs,ms}.

\section{Completion Estimate}
38\%

This is a post hoc, heuristic estimate for the full catalogue target.
A local expenditure identity separates consumption, factor-income, asset and recycling channels under fixed-endowment accounting and avoids duplicate capitalization. Full dynamic incidence and credible channel identification remain unresolved, while unions of pointwise intervals do not establish simultaneous coverage of the complete burden set.

\begin{thebibliography}{9}
\bibitem{bd} S. Borenstein, L. W. Davis, The distributional effects of US clean energy tax credits, \emph{Tax Policy and the Economy} 30 (2016), 191--234.
\bibitem{cfs} J. A. Cronin, D. Fullerton, S. Sexton, Vertical and horizontal redistributions from a carbon tax and rebate, \emph{JAERE} 6(S1) (2019), S169--S208.
\bibitem{ms} G. E. Metcalf, J. H. Stock, The macroeconomic impact of Europe's carbon taxes, \emph{AEJ: Macro} 15(3) (2023), 265--286.
\bibitem{im} G. W. Imbens, C. F. Manski, Confidence intervals for partially identified parameters, \emph{Econometrica} 72 (2004), 1845--1857.
\end{thebibliography}
\end{document}
